% ============================================================
\chapter{Database Normalization}
\label{ch:normalization}
% ============================================================

Why does a well-designed database never contain redundancy? The answer lies in a single word: \emph{anomalies}. When the same information is stored in multiple places, a partial update creates inconsistencies; a deletion can destroy data one intended to keep; an insertion may be impossible without inventing fictitious values. Edgar Codd, as early as 1972, formalized these problems by introducing \emph{functional dependencies} and \emph{normal forms}. The idea is systematic: decompose tables to eliminate redundancy while preserving information (lossless decomposition) and constraints (dependency preservation). First normal form (1NF), second (2NF), third (3NF), and Boyce--Codd normal form (BCNF) form an increasingly strict hierarchy. Understanding this hierarchy is understanding the art of designing robust relational schemas.

\section{Design Problems and Anomalies}

A poorly designed schema leads to \emph{anomalies}:

\begin{definition}[Data Anomalies]
\begin{itemize}
    \item \textbf{Insertion anomaly}: inability to insert data without also
          inserting other unwanted data.
    \item \textbf{Update anomaly}: a modification must be repeated across
          multiple rows to maintain consistency.
    \item \textbf{Deletion anomaly}: deleting one piece of information causes
          the unintended loss of other information.
\end{itemize}
\end{definition}

\begin{example}
Consider the following unnormalized table:

\begin{center}
\begin{tabular}{llllll}
\toprule
\textbf{s\_id} & \textbf{name} & \textbf{course} & \textbf{prof} & \textbf{dept} & \textbf{grade} \\
\midrule
1 & Brown & DB & Smith & CS & 88.5 \\
1 & Brown & Algo & Johnson & Math & 76.0 \\
2 & Davis & DB & Smith & CS & 72.0 \\
\bottomrule
\end{tabular}
\end{center}

\begin{itemize}
    \item \textbf{Update}: if Smith changes department, \emph{all} rows where
          Smith appears must be modified.
    \item \textbf{Insertion}: we cannot record a new course without enrolling a student.
    \item \textbf{Deletion}: deleting Davis's DB enrollment may lose course info
          if it is the last row.
\end{itemize}
\end{example}

\section{Functional Dependencies}

\begin{definition}[Functional Dependency]
Let $R$ be a relation and $X, Y \subseteq \text{attr}(R)$. We say $Y$
\emph{functionally depends} on $X$, denoted $X \to Y$, if for all pairs
of tuples $t_1, t_2 \in R$:
\[
t_1[X] = t_2[X] \implies t_1[Y] = t_2[Y]
\]
In other words, the value of $X$ uniquely determines the value of $Y$.
\end{definition}

\begin{example}
In the University database:
\begin{itemize}
    \item $\text{student\_id} \to \text{last\_name}, \text{first\_name}, \text{major}$
    \item $\text{course\_id} \to \text{title}, \text{credits}, \text{prof\_id}$
    \item $\text{prof\_id} \to \text{prof\_name}, \text{department}$
    \item $(\text{student\_id}, \text{course\_id}) \to \text{grade}$
\end{itemize}
\end{example}

\begin{definition}[Trivial Functional Dependency]
$X \to Y$ is \emph{trivial} if $Y \subseteq X$.
\end{definition}

\begin{definition}[Partial and Transitive Dependencies]
\begin{itemize}
    \item \textbf{Partial}: $X \to Y$ is partial if there exists $X' \subsetneq X$
          such that $X' \to Y$.
    \item \textbf{Transitive}: if $X \to Z$ and $Z \to Y$ (where $Z$ is not a key),
          then $X \to Y$ is transitive.
\end{itemize}
\end{definition}

\section{Armstrong's Axioms}

\begin{theorem}[Armstrong's Axioms]
Armstrong's axioms are \emph{sound} and \emph{complete}:
\begin{enumerate}
    \item \textbf{Reflexivity}: if $Y \subseteq X$, then $X \to Y$.
    \item \textbf{Augmentation}: if $X \to Y$, then $XZ \to YZ$.
    \item \textbf{Transitivity}: if $X \to Y$ and $Y \to Z$, then $X \to Z$.
\end{enumerate}
Derived rules:
\begin{itemize}
    \item \textbf{Union}: if $X \to Y$ and $X \to Z$, then $X \to YZ$.
    \item \textbf{Decomposition}: if $X \to YZ$, then $X \to Y$ and $X \to Z$.
    \item \textbf{Pseudotransitivity}: if $X \to Y$ and $WY \to Z$, then $WX \to Z$.
\end{itemize}
\end{theorem}

\section{Attribute Closure}

\begin{definition}[Closure $X^+$]
The \emph{closure} of $X$ with respect to a set of FDs $F$, denoted $X^+_F$,
is the set of all attributes $A$ such that $X \to A$ can be derived from $F$
using Armstrong's axioms.
\end{definition}

\begin{proposition}[Closure Algorithm]
To compute $X^+$:
\begin{enumerate}
    \item Initialize $X^+ = X$.
    \item Repeat: for each FD $Y \to Z$ in $F$, if $Y \subseteq X^+$,
          add $Z$ to $X^+$.
    \item Stop when $X^+$ no longer changes.
\end{enumerate}
\end{proposition}

\begin{example}
Let $R(A, B, C, D, E)$ with $F = \{A \to B, \; BC \to D, \; D \to E\}$.

Compute $\{A, C\}^+$:
\begin{enumerate}
    \item $X^+ = \{A, C\}$
    \item $A \to B$ and $A \in X^+$: $X^+ = \{A, B, C\}$
    \item $BC \to D$ and $\{B,C\} \subseteq X^+$: $X^+ = \{A, B, C, D\}$
    \item $D \to E$ and $D \in X^+$: $X^+ = \{A, B, C, D, E\}$
\end{enumerate}
Since $X^+$ contains all attributes, $\{A, C\}$ is a candidate key.
\end{example}

\section{First Normal Form (1NF)}

\begin{definition}[1NF]
A relation is in \emph{first normal form} (1NF) if all its attributes are
\emph{atomic}: each cell contains a single, indivisible value.
\end{definition}

\begin{example}
Table NOT in 1NF:
\begin{center}
\begin{tabular}{lll}
\toprule
id & name & phones \\
\midrule
1 & Brown & 555-1234, 555-5678 \\
\bottomrule
\end{tabular}
\end{center}

Table in 1NF (after decomposition):
\begin{center}
\begin{tabular}{ll}
\toprule
id & name \\
\midrule
1 & Brown \\
\bottomrule
\end{tabular}
\qquad
\begin{tabular}{ll}
\toprule
id & phone \\
\midrule
1 & 555-1234 \\
1 & 555-5678 \\
\bottomrule
\end{tabular}
\end{center}
\end{example}

\section{Second Normal Form (2NF)}

\begin{definition}[2NF]
A relation is in \emph{second normal form} if it is in 1NF and every
non-key attribute \emph{fully} depends on the primary key (no partial
functional dependency).
\end{definition}

\begin{remark}
2NF is only relevant for relations with a composite key. If the key is a
single attribute, a relation in 1NF is automatically in 2NF.
\end{remark}

\begin{example}
Let $R(\underline{\text{s\_id}}, \underline{\text{c\_id}}, \text{s\_name}, \text{grade})$
with key $(\text{s\_id}, \text{c\_id})$.

$\text{s\_id} \to \text{s\_name}$ is a partial FD (depends on only part of
the key). Not in 2NF.

Decomposition:
\begin{itemize}
    \item $R_1(\underline{\text{s\_id}}, \text{s\_name})$
    \item $R_2(\underline{\text{s\_id}}, \underline{\text{c\_id}}, \text{grade})$
\end{itemize}
\end{example}

\section{Third Normal Form (3NF)}

\begin{definition}[3NF]
A relation is in \emph{third normal form} if it is in 2NF and no non-key
attribute transitively depends on the primary key.

Formally, for every non-trivial FD $X \to A$:
\begin{itemize}
    \item either $X$ contains a candidate key (superkey),
    \item or $A$ is a prime attribute (belongs to some candidate key).
\end{itemize}
\end{definition}

\begin{example}
Let $R(\underline{\text{s\_id}}, \text{major}, \text{major\_advisor})$.

FDs: $\text{s\_id} \to \text{major}$ and $\text{major} \to \text{major\_advisor}$.

The FD $\text{s\_id} \to \text{major\_advisor}$ is transitive. Not in 3NF.

Decomposition:
\begin{itemize}
    \item $R_1(\underline{\text{s\_id}}, \text{major})$
    \item $R_2(\underline{\text{major}}, \text{major\_advisor})$
\end{itemize}
\end{example}

\section{Boyce-Codd Normal Form (BCNF)}

\begin{definition}[BCNF]
A relation is in \emph{BCNF} if, for every non-trivial functional dependency
$X \to Y$, $X$ is a superkey.
\end{definition}

\begin{theorem}[BCNF vs.\ 3NF]
\begin{itemize}
    \item BCNF $\Rightarrow$ 3NF $\Rightarrow$ 2NF $\Rightarrow$ 1NF.
    \item BCNF is stricter than 3NF: it forbids even dependencies where the
          determinant is not a superkey, even if the dependent is a prime attribute.
    \item Decomposition into BCNF does not always preserve FDs.
    \item Decomposition into 3NF always preserves FDs.
\end{itemize}
\end{theorem}

\begin{example}
Let $R(\underline{A}, \underline{B}, C)$ with FDs: $AB \to C$ and $C \to A$.

$C \to A$: $C$ is not a superkey (the keys are $\{A, B\}$ and $\{B, C\}$).
Not in BCNF.

But $A$ is prime (part of key $\{A, B\}$). So it is in 3NF.
\end{example}

\section{Lossless Decomposition}

\begin{theorem}[Lossless Join Decomposition]
A decomposition of $R$ into $R_1$ and $R_2$ is \emph{lossless} if and only
if one of the following FDs holds:
\[
R_1 \cap R_2 \to R_1 - R_2 \quad \text{or} \quad R_1 \cap R_2 \to R_2 - R_1
\]
i.e., the intersection of the schemas determines at least one of the complements.
\end{theorem}

\section{BCNF Decomposition Algorithm}

\begin{proposition}[BCNF Algorithm]
\begin{enumerate}
    \item Check if the relation is in BCNF.
    \item If not, find an FD $X \to Y$ violating BCNF (where $X$ is not a superkey).
    \item Decompose into:
          \begin{itemize}
              \item $R_1 = X \cup Y$ (with key $X$)
              \item $R_2 = R - Y \cup X$ (all attributes except $Y$, but including $X$)
          \end{itemize}
    \item Repeat on $R_1$ and $R_2$ if necessary.
\end{enumerate}
\end{proposition}

\section{Normal Forms Summary}

\begin{keyformulas}[title=\textbf{Normal Forms Overview}]
\begin{center}
\begin{tabular}{lp{10cm}}
\toprule
\textbf{NF} & \textbf{Condition} \\
\midrule
1NF & All attributes are atomic \\
2NF & 1NF + no partial FD of a non-key attribute on the key \\
3NF & 2NF + no transitive FD of a non-key attribute on the key \\
BCNF & For every non-trivial FD $X \to Y$, $X$ is a superkey \\
\bottomrule
\end{tabular}
\end{center}
\end{keyformulas}

\section{3NF Synthesis (Bernstein's Algorithm)}

\begin{proposition}[3NF Synthesis Algorithm]
\begin{enumerate}
    \item Compute a minimal cover $F_{\min}$ of the FDs.
    \item For each FD $X \to Y$ in $F_{\min}$, create a relation $R_i(X, Y)$.
    \item If no relation contains a candidate key of $R$, add a relation
          with a candidate key.
    \item Eliminate redundant relations (those included in another).
\end{enumerate}
This algorithm guarantees:
\begin{itemize}
    \item Lossless decomposition.
    \item Preservation of functional dependencies.
    \item Result in 3NF.
\end{itemize}
\end{proposition}

\section{Complete Normalization Example}

\begin{codeblock}[title=\textbf{Unnormalized Table}]
\begin{minted}{sql}
CREATE TABLE Enrollments_Denorm (
    student_id INTEGER,
    student_name TEXT,
    major TEXT,
    major_advisor TEXT,
    course_id INTEGER,
    course_title TEXT,
    prof_id INTEGER,
    prof_name TEXT,
    grade REAL
);
\end{minted}
\end{codeblock}

Identified FDs:
\begin{itemize}
    \item $\text{student\_id} \to \text{student\_name}, \text{major}$
    \item $\text{major} \to \text{major\_advisor}$
    \item $\text{course\_id} \to \text{course\_title}, \text{prof\_id}$
    \item $\text{prof\_id} \to \text{prof\_name}$
    \item $(\text{student\_id}, \text{course\_id}) \to \text{grade}$
\end{itemize}

\begin{codeblock}[title=\textbf{Result After 3NF/BCNF Normalization}]
\begin{minted}{sql}
CREATE TABLE Students (
    student_id INTEGER PRIMARY KEY,
    last_name TEXT NOT NULL,
    major TEXT REFERENCES Majors(major)
);

CREATE TABLE Majors (
    major TEXT PRIMARY KEY,
    advisor TEXT NOT NULL
);

CREATE TABLE Professors (
    prof_id INTEGER PRIMARY KEY,
    last_name TEXT NOT NULL
);

CREATE TABLE Courses (
    course_id INTEGER PRIMARY KEY,
    title TEXT NOT NULL,
    prof_id INTEGER REFERENCES Professors(prof_id)
);

CREATE TABLE Enrollments (
    student_id INTEGER REFERENCES Students(student_id),
    course_id INTEGER REFERENCES Courses(course_id),
    grade REAL,
    PRIMARY KEY (student_id, course_id)
);
\end{minted}
\end{codeblock}

\section{Controlled Denormalization}

\begin{remark}
Perfect normalization is not always the best choice in practice.
\emph{Denormalization} consists of intentionally reintroducing redundancy
to improve read performance (avoiding costly joins).
\end{remark}

\begin{bestpractice}[title=\textbf{When to Denormalize?}]
\begin{itemize}
    \item When reads are much more frequent than writes.
    \item For reporting tables or data warehouses.
    \item Always document and justify denormalization.
    \item Use triggers or materialized views to maintain consistency.
\end{itemize}
\end{bestpractice}

\section*{Exercises}

\begin{exercise}
Let $R(A, B, C, D, E)$ with $F = \{A \to BC, \; C \to D, \; BD \to E\}$.
\begin{enumerate}
    \item Compute $\{A\}^+$.
    \item Identify the candidate keys.
    \item Is the relation in 2NF? 3NF? BCNF? Justify.
    \item Decompose into BCNF.
\end{enumerate}
\end{exercise}

\begin{exercise}
Normalize the following relation into 3NF using the synthesis algorithm:
$R(\text{OrderNum}, \text{OrderDate}, \text{CustNum}, \text{CustName},
\text{ProdNum}, \text{ProdName}, \text{Quantity}, \text{UnitPrice})$

FDs: $\text{OrderNum} \to \text{OrderDate}, \text{CustNum}$;
$\text{CustNum} \to \text{CustName}$;
$\text{ProdNum} \to \text{ProdName}, \text{UnitPrice}$;
$(\text{OrderNum}, \text{ProdNum}) \to \text{Quantity}$.
\end{exercise}

\begin{exercise}
Give an example of a relation that is in 3NF but not in BCNF. Decompose
it into BCNF and verify whether the functional dependencies are preserved.
\end{exercise}

\begin{exercise}
Verify that the decomposition of $R(A, B, C)$ with $F = \{A \to B\}$ into
$R_1(A, B)$ and $R_2(A, C)$ is lossless. Use Heath's theorem.
\end{exercise}
